\subsection{Hierarchy of Quark Masses, Cabibbo Angles and CP Violation}

\textbf{Authors:} C. D. Froggatt, Holger Bech Nielsen

\paragraph{主な主張}
基礎的な結合が同程度でも、近似的に保存される量子数の違いによってクォーク質量に大きな階層が生じ、混合角と質量比の大きさに関係が現れ得ると提案した\cite{Froggatt:1978nt}。

\paragraph{新規性}
小さな対称性破れの因子を電荷に応じた次数で有効結合へ導入することで、Yukawa結合の階層を説明する機構を提示した。

\paragraph{位置付け}
フレーバー対称性によるフェルミオン質量階層の説明に広く用いられるFroggatt--Nielsen機構の原論文である。

\paragraph{コメント（読書メモ）}
読書メモ：Froggatt--Nielsen模型の原論文。
